% ---
% Conclusão
% ---
\chapter{Conclusão}

Este projeto de TCC propõe o desenvolvimento e a avaliação de um fluxo padronizado para extração estruturada de informações em provas objetivas manuscritas. O foco está na comparação entre dois modelos multimodais acessados por APIs hospedadas: um modelo aberto e um modelo proprietário. O aplicativo móvel será utilizado apenas para captura das imagens e visualização dos resultados, enquanto o processamento será realizado pelo backend e pelas APIs dos modelos avaliados.

A expectativa de conclusão é identificar, com base no experimento planejado, qual abordagem apresenta melhor equilíbrio entre acurácia, validade da saída JSON, latência, custo operacional, robustez e necessidade de validação humana. Assim, o resultado esperado não é apenas demonstrar um aplicativo funcionando, mas produzir uma análise comparativa útil para decisões técnicas em Sistemas de Informação.

O estudo permanecerá limitado ao conjunto de provas, modelos, APIs, prompts e métricas definidos no projeto. Mesmo assim, espera-se que a avaliação permita discutir de forma objetiva os trade-offs entre um modelo multimodal aberto e um modelo multimodal proprietário ambos acessados via API.
